\subsection{Coefficient geometry and the regime decision map}
\label{sec:coefficient-regimes}

Different relative sizes of the five coefficients in \eqref{eq:slice}
lead to different passage scales. We introduce these scales here and
summarize the resulting cases in Figure~\ref{fig:coefficient-phase-decision}.
The corresponding bounds are proved in the later face propositions.
Before counting zeros, every phase chart is transported to the physical
sections of \eqref{eq:return}.

Translate the zero of the transverse eigenvalue to the origin and put
\begin{equation}
 \xi=v-\gamma,
 \qquad \alpha _0=P(\gamma),
 \qquad \alpha _1=P'(\gamma),
 \qquad q_\gamma=1+2\gamma,
 \qquad \ell(v)=1-Av+Cv^2.
 \label{deg:central-coordinates}
\end{equation}
Then \eqref{eq:raw} becomes
\begin{equation}
 \begin{split}
  \dot \xi&=\alpha _0+\alpha _1\xi-3\gamma \xi^2-\xi^3
                 -\ell(\gamma+\xi)z,\\
  \dot z&=-z\{q_\gamma \xi+\xi^2+C(\gamma+\xi)z\}.
 \end{split}
 \label{deg:central-field}
\end{equation}
Thus \((A,C,\gamma,\alpha _0,\alpha _1)\) are analytic coordinates on
the five-parameter slice near the organizing field. After one common
coefficient shrink, \(q_\gamma\) and \(\ell\) are positive units on the
fixed raw rectangle, and the regular transition maps between fixed sections
and their inverses have uniform bounds through every derivative order fixed
in advance.

For negative opening write
\begin{equation}
 \alpha _0=-s^2,
 \qquad p=\frac{\alpha _1}{s},
 \qquad v=\gamma+sX,
 \qquad z=s^2Y.
 \label{deg:source-dictionary}
\end{equation}
This gives the source field \eqref{ann:source-field} with parameters
\((a,p,\gamma,s)=(-1,p,\gamma,s)\). Two exact compactifications cover
its small-\(p\) coefficient ends. For \(\gamma\ne0\), set
\begin{equation}
 g=|\gamma|,\qquad k=\frac{s}{g^{3/2}},\qquad
 \pi=\frac{p}{\sqrt g},\qquad \mathfrak q=k\pi .
 \label{deg:gamma-chart}
\end{equation}
Here \(g\) is always a coefficient radius, never a section height. In the
cubic chart set
\begin{equation}
 t=s^{1/3},\qquad U=\frac{p}{t},\qquad
 G=\frac{\gamma}{t^2},\qquad \tau_{\rm cub}=t^2.
 \label{deg:cubic-chart}
\end{equation}
We write \(\tau_{\rm cub}\) for this cubic scale throughout. The
far-root, two-saddle and degree-loss arguments each use an unadorned
\(\tau\) for their local scale; any relation such as \(\tau=p^2\)
is stated where that scale is introduced.

\begin{figure}[H]
\centering
\small
\fbox{\begin{minipage}{0.94\linewidth}
\centering
\(\theta\) in the five-parameter slice
\(\longrightarrow (\gamma,\alpha _0,\alpha _1;A,C)\)

\medskip
\begin{tabular}{@{}p{.27\linewidth}@{\(\longrightarrow\)}p{.62\linewidth}@{}}
\shortstack[l]{\(K\le0\) or no\\through passage}
 & stopped/no-passage endpoint; it is not a periodic-orbit zero\\[2pt]
\(K>0\), high central height
 & affine high-height estimate, Corollary~\ref{rew:cor-high-derivative-signed}\\[2pt]
\(\alpha _0\ge0\), low height
 & determinant chart or signed trace chart,
   Lemma~\ref{deg:lem-nonnegative-domain}\\[2pt]
\(\alpha _0<0,\ p\le0\)
 & raw Jacobian contraction, Lemma~\ref{deg:lem-raw-jacobian}\\[2pt]
\(\alpha _0<0,\ 0<p\ll1\)
 & gamma-dominant, cubic, degree-loss, or remote-root face,
   Lemma~\ref{deg:lem-small-source-routing}\\[2pt]
\(\alpha _0<0,\ p\asymp1\)
 & compact positive-ratio passage,
   Lemma~\ref{ann:lem-compact-positive-domains}\\[2pt]
\(\alpha _0<0,\ p\gg1\)
 & signed-trace terminal block plus compact annuli,
   Lemma~\ref{deg:lem-large-source-routing}
\end{tabular}

\medskip
Each row \(\longrightarrow\) zeros of the same exact displacement
\(x-m\Pi(x)\) on one coefficient fibre
\(\longrightarrow\) one bound on each whole phase fibre, then the
maximum across coefficient alternatives
\(\longrightarrow\) one common coefficient neighborhood and one common
Hausdorff collar
\(\longrightarrow\) lift from five to twelve coefficients.
\end{minipage}}
\caption{Global coefficient--phase decision map. The arrows specify the
logical routing, not additional coordinate changes. Strict overlaps between
adjacent phase and coefficient blocks are recorded in
Table~\ref{tab:coefficient-phase-cover}.}
\label{fig:coefficient-phase-decision}
\end{figure}

For later reference, the symbols with global scope are
\(\theta=(A,C,\mu_1,\mu_2,\mu_3)\),
\((\gamma,\alpha _0,\alpha _1)\), the finite multiplier \(m\), and the
actual input of \eqref{eq:return}. Symbols such as \(g,k,\pi,\mathfrak q\),
\(t,U,G\), and the degree-loss exponents are local to the chart in which
they are defined. Actual phase heights carry labels such as \(h\),
\(h_L\), \(V\), \(\widehat V\), \(W\), or \(\Theta\); they are never
identified with a coefficient radius or with a stopped trace gap.
